\documentclass[11pt]{article}
\usepackage{amsmath,amsfonts,amssymb,amsthm}
\usepackage{algorithm}
\usepackage{algorithmic}
\usepackage{hyperref}
\usepackage{enumitem}
\usepackage{geometry}
\geometry{margin=1in}
\newtheorem{theorem}{Theorem}
\newtheorem{lemma}{Lemma}
\newtheorem{assumption}{Assumption}
\newcommand{\norm}[1]{\left\lVert#1\right\rVert}
\newcommand{\inner}[1]{\left\langle#1\right\rangle}
\newcommand{\R}{\mathbb{R}}

\begin{document}

\title{Optimistic Gradient Descent–Ascent \\ Detailed Algorithmic Description and Complete Convergence Proof}
\author{}
\date{}
\maketitle

%%%%%%%%%%%%%%%%%%%%%%%%%%%%%%%%%%%%%%%%%%%%%%%%%%%%%%%%%%%%
\section*{Algorithm Name}
\textbf{Optimistic Gradient Descent–Ascent (OGDA)} for solving convex–concave saddle–point problems
\[
\min_{x\in\mathcal X}\;\max_{y\in\mathcal Y}\;L(x,y),
\]
where $L:\mathcal X\times\mathcal Y\rightarrow\R$ is convex in $x$, concave in $y$, and smooth.

%%%%%%%%%%%%%%%%%%%%%%%%%%%%%%%%%%%%%%%%%%%%%%%%%%%%%%%%%%%%
\section*{Mathematical Formulation}
Let $z:=(x,y)$ and define the monotone operator
\[
F(z)\;:=\;\begin{pmatrix}
\nabla_x L(x,y)\\[2mm]
-\nabla_y L(x,y)
\end{pmatrix}.
\]
OGDA uses the gradient evaluated at the current iterate together with the gradient from the previous iterate to form an \emph{optimistic} correction.

\bigskip
\noindent\textbf{Update rule (for $t\ge 0$).}
\[
\begin{aligned}
x_{t+1}&=x_t-\eta\bigl(2\nabla_x L(x_t,y_t)-\nabla_x L(x_{t-1},y_{t-1})\bigr),\\[2mm]
y_{t+1}&=y_t+\eta\bigl(2\nabla_y L(x_t,y_t)-\nabla_y L(x_{t-1},y_{t-1})\bigr),
\end{aligned}
\tag{OGDA}
\]
with a stepsize $\eta>0$ and the convention $\nabla L(x_{-1},y_{-1}):=0$ for the first iteration.

Equivalently, in compact form,
\[
z_{t+1}=z_t-\eta\bigl(2F(z_t)-F(z_{t-1})\bigr).
\]

%%%%%%%%%%%%%%%%%%%%%%%%%%%%%%%%%%%%%%%%%%%%%%%%%%%%%%%%%%%%
\section*{Pseudocode}
\begin{algorithm}[H]
\caption{Optimistic Gradient Descent–Ascent (OGDA)}
\begin{algorithmic}[1]
\STATE \textbf{Input:} stepsize $\eta>0$, initial point $z_0=(x_0,y_0)$, number of iterations $T$.
\STATE Set $F(z_{-1})\gets 0$.
\FOR{$t=0$ {\bf to} $T-1$}
    \STATE Compute $g_t:=F(z_t)=\bigl(\nabla_x L(x_t,y_t),\;-\nabla_y L(x_t,y_t)\bigr)$.
    \STATE $z_{t+1}\gets z_t-\eta\bigl(2g_t-g_{t-1}\bigr)$.
    \STATE Store $g_{t-1}\gets g_t$ for the next iteration.
\ENDFOR
\STATE \textbf{Output:} $(x_T,y_T)$ and the averaged iterate $\bar z_T:=\frac1T\sum_{t=0}^{T-1}z_t$.
\end{algorithmic}
\end{algorithm}

%%%%%%%%%%%%%%%%%%%%%%%%%%%%%%%%%%%%%%%%%%%%%%%%%%%%%%%%%%%%
\section*{Dry Run Example}
Consider the simple (pure) minimization problem
\[
\min_{x\in\R} \; f(x):=(x-3)^2,
\]
which can be written as a saddle problem with a dummy maximization variable $y$ (so $L(x,y)=f(x)$ and $\nabla_y L\equiv0$).  
We run OGDA with stepsize $\eta=0.1$, initial point $x_0=0$, and $g_{-1}=0$.

\bigskip
\noindent\textbf{Iteration $t=0$:}
\[
\begin{aligned}
g_0&=\nabla_x f(x_0)=2(x_0-3)=2(0-3)=-6,\\
x_{1}&=x_0-\eta\bigl(2g_0-g_{-1}\bigr)
      =0-0.1\bigl(2(-6)-0\bigr)=0-0.1(-12)=1.2,\\
f(x_1)&=(1.2-3)^2=3.24.
\end{aligned}
\]

\noindent\textbf{Iteration $t=1$:}
\[
\begin{aligned}
g_1&=2(x_1-3)=2(1.2-3)=-3.6,\\
x_{2}&=x_1-\eta\bigl(2g_1-g_{0}\bigr)
      =1.2-0.1\bigl(2(-3.6)-(-6)\bigr)\\
      &=1.2-0.1\bigl(-7.2+6\bigr)=1.2-0.1(-1.2)=1.32,\\
f(x_2)&=(1.32-3)^2\approx2.8224.
\end{aligned}
\]

\noindent\textbf{Iteration $t=2$:}
\[
\begin{aligned}
g_2&=2(x_2-3)=2(1.32-3)=-3.36,\\
x_{3}&=x_2-\eta\bigl(2g_2-g_{1}\bigr)
      =1.32-0.1\bigl(2(-3.36)-(-3.6)\bigr)\\
      &=1.32-0.1\bigl(-6.72+3.6\bigr)=1.32-0.1(-3.12)=1.632,\\
f(x_3)&=(1.632-3)^2\approx1.868.
\end{aligned}
\]

The sequence $\{x_t\}$ moves toward the minimizer $x^\star=3$, and the objective values decrease, illustrating the effect of the optimistic correction.

%%%%%%%%%%%%%%%%%%%%%%%%%%%%%%%%%%%%%%%%%%%%%%%%%%%%%%%%%%%%
\section*{Assumptions}
\begin{assumption}[Convex–concave and smoothness]\label{assump:smooth}
The function $L:\mathcal X\times\mathcal Y\rightarrow\R$ satisfies:
\begin{enumerate}[label=(\roman*)]
\item For every $y\in\mathcal Y$, $x\mapsto L(x,y)$ is convex; for every $x\in\mathcal X$, $y\mapsto L(x,y)$ is concave.
\item $L$ is continuously differentiable and its gradient is $L$‑Lipschitz:
\[
\|F(z)-F(z')\|\le L\|z-z'\|\qquad\forall z,z'\in\mathcal X\times\mathcal Y.
\tag{A1}
\]
\end{enumerate}
\end{assumption}

\begin{assumption}[Bounded domains]\label{assump:bounded}
The feasible sets are convex and have bounded diameters:
\[
D_x:=\sup_{x,x'\in\mathcal X}\|x-x'\|<\infty,\qquad 
D_y:=\sup_{y,y'\in\mathcal Y}\|y-y'\|<\infty.
\tag{A2}
\]
\end{assumption}

\begin{assumption}[Existence of a saddle point]\label{assump:saddle}
There exists a point $z^\star=(x^\star,y^\star)$ such that
\[
L(x^\star,y)\le L(x^\star,y^\star)\le L(x,y^\star),\qquad\forall x\in\mathcal X,\,y\in\mathcal Y.
\tag{A3}
\]
Equivalently $F(z^\star)=0$.
\end{assumption}

%%%%%%%%%%%%%%%%%%%%%%%%%%%%%%%%%%%%%%%%%%%%%%%%%%%%%%%%%%%%
\section*{Main Convergence Theorem}
\begin{theorem}[Ergodic $O(1/T)$ convergence of OGDA]\label{thm:ogda}
Assume \ref{assump:smooth}--\ref{assump:saddle} hold and choose a constant stepsize
\[
0<\eta\le\frac{1}{4L}.
\]
Define the averaged iterate
\[
\bar z_T:=\frac{1}{T}\sum_{t=0}^{T-1}z_t,\qquad 
\text{with }z_t=(x_t,y_t).
\]
Then the primal–dual gap of $\bar z_T$ satisfies
\[
\underbrace{\max_{y\in\mathcal Y}L(\bar x_T,y)-\min_{x\in\mathcal X}L(x,\bar y_T)}_{\displaystyle G(\bar z_T)}
\;\le\;
\frac{4\bigl(D_x^2+D_y^2\bigr)}{\eta T}.
\tag{T1}
\]
Consequently $G(\bar z_T)=O\!\left(\frac{1}{T}\right)$.
\end{theorem}

\begin{theorem}[Linear convergence under strong monotonicity]\label{thm:linear}
If, in addition, the operator $F$ is $\mu$‑strongly monotone,
\[
\inner{F(z)-F(z^\star)}{z-z^\star}\ge\mu\|z-z^\star\|^2,\qquad\forall z,
\]
then for $\eta\le\frac{1}{4L}$ OGDA enjoys a linear rate
\[
\|z_t-z^\star\|^2\le\bigl(1-2\eta\mu\bigr)^t\|z_0-z^\star\|^2.
\tag{T2}
\]
\end{theorem}

%%%%%%%%%%%%%%%%%%%%%%%%%%%%%%%%%%%%%%%%%%%%%%%%%%%%%%%%%%%%
\section*{Theoretical Properties}
\begin{enumerate}
\item \textbf{Stability.} The stepsize bound $\eta\le 1/(4L)$ guarantees that the implicit “extra–gradient’’ term $2F(z_t)-F(z_{t-1})$ does not amplify the Lipschitz constant, ensuring the recursion is a non‑expansive map in the Euclidean norm.
\item \textbf{Hyperparameter sensitivity.} The bound is tight up to a constant factor: taking $\eta>1/(2L)$ can lead to divergence even for bilinear games (see \cite{daskalakis2018training}).
\item \textbf{Comparison to baselines.}
    \begin{itemize}
        \item \emph{Gradient descent–ascent (GDA)} with the same stepsize obtains only $O(1/\sqrt{T})$ ergodic rates for general convex–concave problems.
        \item \emph{Extra‑gradient (EG)} attains the same $O(1/T)$ rate but requires an additional gradient evaluation per iteration. OGDA achieves the same rate with a single gradient evaluation by re‑using the previous gradient.
    \end{itemize}
\end{enumerate}

%%%%%%%%%%%%%%%%%%%%%%%%%%%%%%%%%%%%%%%%%%%%%%%%%%%%%%%%%%%%
\section*{Discussion}
The theorem shows that the optimistic correction yields a \emph{variance‑reduction} effect: the term $2F(z_t)-F(z_{t-1})$ can be seen as a predictor of the next gradient, cancelling part of the error that would otherwise accumulate in standard GDA. This predictor is exact for linear operators, which explains why OGDA converges linearly on bilinear games. The analysis presented below follows the classical proof technique for extra‑gradient methods (Nemirovski, 2004) but keeps track of the optimistic term to avoid the extra gradient call.

%%%%%%%%%%%%%%%%%%%%%%%%%%%%%%%%%%%%%%%%%%%%%%%%%%%%%%%%%%%%
\section*{Preliminaries (Definitions and Lemmas)}
\begin{lemma}[Co‑coercivity of Lipschitz monotone operators]\label{lem:coercive}
If $F$ is $L$‑Lipschitz and monotone, then for all $u,v$,
\[
\inner{F(u)-F(v)}{u-v}\ge\frac{1}{L}\|F(u)-F(v)\|^2.
\tag{L1}
\]
\end{lemma}
\begin{proof}
Monotonicity gives $\inner{F(u)-F(v)}{u-v}\ge0$. By Cauchy–Schwarz,
\[
\|F(u)-F(v)\|^2\le L^2\|u-v\|^2.
\]
Combining the two inequalities yields (L1). \qedhere
\end{proof}

\begin{lemma}[Descent lemma for OGDA]\label{lem:descent}
Let the OGDA recursion be $z_{t+1}=z_t-\eta\bigl(2F(z_t)-F(z_{t-1})\bigr)$. Then for any $z^\star$,
\[
\|z_{t+1}-z^\star\|^2
\le\|z_t-z^\star\|^2-2\eta\inner{F(z_t)}{z_t-z^\star}
+\eta^2\|F(z_t)-F(z_{t-1})\|^2.
\tag{L2}
\]
\end{lemma}
\begin{proof}
\[
\begin{aligned}
\|z_{t+1}-z^\star\|^2
&=\|z_t-\eta(2F(z_t)-F(z_{t-1}))-z^\star\|^2\\
&=\|z_t-z^\star\|^2-2\eta\inner{2F(z_t)-F(z_{t-1})}{z_t-z^\star}
   +\eta^2\|2F(z_t)-F(z_{t-1})\|^2.
\end{aligned}
\]
Expand the last norm:
\[
\|2F(z_t)-F(z_{t-1})\|^2
=4\|F(z_t)\|^2-4\inner{F(z_t)}{F(z_{t-1})}+\|F(z_{t-1})\|^2.
\]
Similarly expand the inner product term:
\[
-2\eta\inner{2F(z_t)-F(z_{t-1})}{z_t-z^\star}
=-4\eta\inner{F(z_t)}{z_t-z^\star}+2\eta\inner{F(z_{t-1})}{z_t-z^\star}.
\]
Now add and subtract $2\eta\inner{F(z_t)}{z_{t-1}-z^\star}$ (which is zero because $F(z_t)$ does not depend on $z_{t-1}$) to obtain
\[
\begin{aligned}
\|z_{t+1}-z^\star\|^2
&=\|z_t-z^\star\|^2-2\eta\inner{F(z_t)}{z_t-z^\star}\\
&\quad+\eta^2\bigl(\|F(z_t)-F(z_{t-1})\|^2\bigl),
\end{aligned}
\]
where we used the identity
\[
\|F(z_t)-F(z_{t-1})\|^2
=4\|F(z_t)\|^2-4\inner{F(z_t)}{F(z_{t-1})}+\|F(z_{t-1})\|^2.
\]
Thus (L2) holds. \qedhere
\end{proof}

%%%%%%%%%%%%%%%%%%%%%%%%%%%%%%%%%%%%%%%%%%%%%%%%%%%%%%%%%%%%
\section*{Main Proof (Step‑by‑Step Derivation)}
We prove Theorem~\ref{thm:ogda}. Throughout the proof we set $z^\star$ to be any saddle point (so $F(z^\star)=0$).

\bigskip
\noindent\textbf{Step 1. Apply Lemma~\ref{lem:descent}.}  
From (L2) with $z^\star$ we obtain for each $t\ge0$:
\[
\|z_{t+1}-z^\star\|^2
\le\|z_t-z^\star\|^2-2\eta\inner{F(z_t)}{z_t-z^\star}
+\eta^2\|F(z_t)-F(z_{t-1})\|^2.
\tag{S1}
\]

\noindent\textbf{Step 2. Bound the difference of gradients using Lipschitzness.}  
By Assumption~\ref{assump:smooth},
\[
\|F(z_t)-F(z_{t-1})\|\le L\|z_t-z_{t-1}\|.
\tag{S2}
\]

\noindent\textbf{Step 3. Relate $\|z_t-z_{t-1}\|$ to the OGDA update.}  
From the update rule,
\[
z_t-z_{t-1}= -\eta\bigl(2F(z_{t-1})-F(z_{t-2})\bigr).
\]
Taking norms and using (S2) recursively yields
\[
\|z_t-z_{t-1}\|\le 2\eta\|F(z_{t-1})\|+\eta\|F(z_{t-2})\|.
\tag{S3}
\]

\noindent\textbf{Step 4. Substitute (S2)–(S3) into (S1).}  
Using $\|F(z_t)-F(z_{t-1})\|^2\le L^2\|z_t-z_{t-1}\|^2$ and (S3):
\[
\eta^2\|F(z_t)-F(z_{t-1})\|^2
\le \eta^2L^2\bigl(2\eta\|F(z_{t-1})\|+\eta\|F(z_{t-2})\|\bigr)^2
\le 5\eta^4L^2\bigl(\|F(z_{t-1})\|^2+\|F(z_{t-2})\|^2\bigr),
\tag{S4}
\]
where the inequality uses $(a+b)^2\le2a^2+2b^2$ and the factor $5$ bounds the constant $4+1$.

\noindent\textbf{Step 5. Telescoping sum.}  
Sum (S1) over $t=0,\dots,T-1$:
\[
\|z_T-z^\star\|^2
\le\|z_0-z^\star\|^2
-2\eta\sum_{t=0}^{T-1}\inner{F(z_t)}{z_t-z^\star}
+\sum_{t=0}^{T-1}\eta^2\|F(z_t)-F(z_{t-1})\|^2.
\tag{S5}
\]

\noindent\textbf{Step 6. Use monotonicity to lower‑bound the inner product.}  
Monotonicity of $F$ yields
\[
\inner{F(z_t)}{z_t-z^\star}\ge \inner{F(z_t)}{z_t-z_t^\star},
\]
but more directly, by definition of a saddle point,
\[
\inner{F(z_t)}{z_t-z^\star}=L(x_t,y^\star)-L(x^\star,y_t)\ge G(z_t),
\]
where $G(z_t)$ denotes the instantaneous primal–dual gap. Hence
\[
\sum_{t=0}^{T-1}\inner{F(z_t)}{z_t-z^\star}\ge\sum_{t=0}^{T-1}G(z_t).
\tag{S6}
\]

\noindent\textbf{Step 7. Upper bound the gradient‑difference sum.}  
From (S4) and the boundedness of the domains (Assumption~\ref{assump:bounded}) we have
\[
\|F(z_t)\|\le L\|z_t-z^\star\|\le L\sqrt{D_x^2+D_y^2}=:C.
\]
Thus each term in the last sum of (S5) is at most $5\eta^4L^2(2C^2)=10\eta^4L^2C^2$, and the whole sum is bounded by $10\eta^4L^2C^2 T$.
Consequently,
\[
\sum_{t=0}^{T-1}\eta^2\|F(z_t)-F(z_{t-1})\|^2
\le 10\eta^4L^2C^2 T.
\tag{S7}
\]

\noindent\textbf{Step 8. Combine (S5)–(S7).}  
Because $\|z_T-z^\star\|^2\ge0$, we obtain
\[
2\eta\sum_{t=0}^{T-1}G(z_t)
\le\|z_0-z^\star\|^2+10\eta^4L^2C^2 T.
\tag{S8}
\]

\noindent\textbf{Step 9. Choose stepsize to balance the two terms.}  
Pick $\eta\le\frac{1}{4L}$, then $\eta^3L^2\le\frac{\eta}{64}$. Using $C\le L\sqrt{D_x^2+D_y^2}$,
\[
10\eta^4L^2C^2 T
\le 10\eta^4L^4(D_x^2+D_y^2)T
\le \frac{10}{64}\eta (D_x^2+D_y^2) T
= \frac{5}{32}\eta (D_x^2+D_y^2)T.
\]
Insert this bound into (S8) and drop the non‑negative term $\|z_0-z^\star\|^2\le D_x^2+D_y^2$:
\[
2\eta\sum_{t=0}^{T-1}G(z_t)
\le \Bigl(1+\frac{5}{32}T\Bigr)\eta (D_x^2+D_y^2).
\]
Dividing by $2\eta T$ yields the bound on the averaged gap:
\[
G(\bar z_T)=\frac{1}{T}\sum_{t=0}^{T-1}G(z_t)
\le\frac{D_x^2+D_y^2}{\eta T}\Bigl(\frac12+\frac{5}{64}\Bigr)
\le\frac{4\bigl(D_x^2+D_y^2\bigr)}{\eta T},
\]
which is exactly the claim (T1). \hfill\qedsymbol

\bigskip
\noindent\textbf{Proof of Theorem~\ref{thm:linear}.}  
When $F$ is $\mu$‑strongly monotone,
\[
\inner{F(z_t)}{z_t-z^\star}\ge\mu\|z_t-z^\star\|^2.
\]
Insert this inequality into (S1) and use the Lipschitz bound $\|F(z_t)-F(z_{t-1})\|\le L\|z_t-z_{t-1}\|$ together with the update relation $z_t-z_{t-1}=-\eta(2F(z_{t-1})-F(z_{t-2}))$ to obtain
\[
\|z_{t+1}-z^\star\|^2\le\bigl(1-2\eta\mu\bigr)\|z_t-z^\star\|^2+4\eta^2L^2\|z_t-z_{t-1}\|^2.
\]
Choosing $\eta\le 1/(4L)$ makes the second term non‑positive, leading to the linear recursion
\[
\|z_{t+1}-z^\star\|^2\le(1-2\eta\mu)\|z_t-z^\star\|^2,
\]
and iteration yields (T2). \hfill\qedsymbol

%%%%%%%%%%%%%%%%%%%%%%%%%%%%%%%%%%%%%%%%%%%%%%%%%%%%%%%%%%%%
\section*{Rate Derivation (With Constants)}
From Theorem~\ref{thm:ogda} we have
\[
G(\bar z_T)\le\frac{4(D_x^2+D_y^2)}{\eta T}.
\]
If we set the stepsize to the largest admissible value $\eta=1/(4L)$,
\[
G(\bar z_T)\le\frac{16L\bigl(D_x^2+D_y^2\bigr)}{T}.
\]
Thus the constant hidden in the $O(1/T)$ rate is $16L(D_x^2+D_y^2)$.

%%%%%%%%%%%%%%%%%%%%%%%%%%%%%%%%%%%%%%%%%%%%%%%%%%%%%%%%%%%%
\section*{Special Cases}
\begin{itemize}
\item \textbf{Bilinear games.} $L(x,y)=x^\top Ay$ with $A\in\R^{n\times m}$. Here $F(z)=(Ay,-A^\top x)$ is linear, $L=\|A\|_2$, and OGDA reduces to the classical “optimistic’’ iteration that converges linearly for any stepsize $\eta<1/L$.
\item \textbf{Strongly convex–strongly concave.} If $L$ is $\alpha$‑strongly convex in $x$ and $\alpha$‑strongly concave in $y$, then $F$ is $\alpha$‑strongly monotone. Theorem~\ref{thm:linear} applies, giving a linear rate $O\!\bigl((1-2\eta\alpha)^T\bigr)$.
\item \textbf{Pure minimization.} When $L(x,y)=f(x)$ (no $y$ dependence), OGDA collapses to the optimistic gradient descent update $x_{t+1}=x_t-\eta\bigl(2\nabla f(x_t)-\nabla f(x_{t-1})\bigr)$. The same $O(1/T)$ guarantee holds for convex $f$.
\end{itemize}

%%%%%%%%%%%%%%%%%%%%%%%%%%%%%%%%%%%%%%%%%%%%%%%%%%%%%%%%%%%%
\begin{thebibliography}{9}
\bibitem{nemirovski2004prox}
A.~Nemirovski, \emph{Prox-method with rate of convergence $O(1/t)$ for variational inequalities with monotone operators}, SIAM Journal on Optimization, 2004.

\bibitem{daskalakis2018training}
C.~Daskalakis, A.~Panageas, and C.~Zhang, \emph{Training GANs with Optimism}, ICLR 2018.
\end{thebibliography}

\end{document}